\section{Pengujian}
\label{sec:pengujian}

Pengujian sistem dilakukan untuk memastikan implementasi memenuhi kebutuhan yang sudah didefinisikan pada Subbab~\ref{subsec:impl-skema-tipe-data} hingga~\ref{subsec:rancangan-publikasi-iota}.

\subsection{Strategi Pengujian}
\label{subsec:strategi-pengujian}

Strategi pengujian disusun berdasarkan kebutuhan sistem pada Tabel~\ref{tab:kebutuhan-fungsional} dan Tabel~\ref{tab:kebutuhan-nonfungsional}. Pengujian dibagi menjadi pengujian kebutuhan fungsional dan pengujian kebutuhan nonfungsional.

Pengujian kebutuhan fungsional diberi ID P-KF dan dijalankan melalui skenario manual untuk memeriksa perilaku sistem dari sudut pandang komponen DecMed sebagai pengirim \textit{audit event} dan ALS sebagai penerima. Pengujian kebutuhan nonfungsional diberi ID P-KNF dan dijalankan melalui pengujian unit atau integrasi terhadap komponen verifikasi tanda tangan digital, mekanisme \textit{hash-chaining}, penulisan berkas log, rotasi dan pengarsipan ke IPFS, serta publikasi metadata ke IOTA. Pengujian otomatis dijalankan dengan \texttt{cargo test} pada komponen yang relevan.

Kriteria keberhasilan pengujian ditentukan berdasarkan kesesuaian hasil skenario manual terhadap kebutuhan fungsional dan kelulusan pengujian otomatis terhadap kebutuhan nonfungsional. Setiap pengujian dikaitkan langsung dengan ID kebutuhan yang diuji untuk memastikan ketertelusuran.

\subsection{Pengujian Kebutuhan Fungsional}
\label{subsec:pengujian-fungsional}

Pengujian kebutuhan fungsional dijalankan secara manual untuk memastikan perilaku sistem sesuai dengan rancangan pada Subbab~\ref{subsec:rancangan-arsitektur} hingga~\ref{subsec:rancangan-publikasi-iota}. Terdapat lima skenario pengujian dengan ID P-KF-01 sampai P-KF-05 yang dijalankan untuk memverifikasi kemampuan ALS menerima, memformat, menyimpan, dan menyajikan kembali \textit{audit event} dari seluruh \textit{event source}.

\subsubsection{P-KF-01: Penerimaan dan Agregasi \textit{Audit Event} dari Berbagai \textit{Event Source}}

Pengujian P-KF-01 memverifikasi KF-TA-01, yaitu kemampuan ALS mengagregasi data log dari berbagai komponen DecMed ke dalam satu repositori terpusat. Skenario pengujian dilakukan dengan mengirimkan \textit{encryptedsignedevent} secara berturut-turut dari empat \textit{event source} yang berbeda, yaitu \textit{Patient Client}, \textit{Hospital Client}, \textit{Ministry Client}, dan PRE Server, ke \textit{endpoint} \texttt{POST /api/events} pada ALS.

Setiap \textit{event source} mengirimkan dua event: satu event dengan hasil \texttt{Success} dan satu event dengan hasil \texttt{Failure}, untuk memastikan ALS dapat menangani seluruh varian \texttt{AuditOutcome}. Setelah seluruh pengiriman selesai, berkas log \texttt{./logs/audit\_logs.log} diperiksa untuk memastikan seluruh delapan event berhasil ditulis sebagai baris JSON-Lines yang terpisah tanpa ada yang terlewat.

Hasil pengujian menunjukkan seluruh delapan \textit{audit record} berhasil tersimpan pada berkas log. Seluruh delapan \texttt{EncryptedSignedEvent} berhasil diterima dan disimpan sebagai \texttt{AuditRecord} yang terpisah. Identitas penanda tangan setiap record dapat ditelusuri melalui \texttt{iota\_address}, sedangkan atribut \texttt{source\_component} beserta atribut event lainnya berada di dalam payload terenkripsi (\textit{ciphertext}) dan baru dapat diperiksa setelah proses dekripsi menggunakan \textit{private key} ALS.

\subsubsection{P-KF-02: Pencatatan Keseluruhan Jenis \textit{Audit Event} EV1--EV13}

Pengujian P-KF-02 memverifikasi KF-TA-02, yaitu kemampuan ALS menangkap dan merekam keseluruhan 13 jenis \textit{audit event} yang telah didefinisikan. Skenario pengujian dilakukan dengan mengirimkan satu \textit{signed event} untuk setiap varian \texttt{AuditEventDetails} (EV1 hingga EV13) ke ALS.

Setiap \textit{event} dikonstruksi dengan nilai bidang yang valid sesuai skema masing-masing varian pada Lampiran~\ref{code:audit-event-details}. Setiap varien EV1-EV13 dibentuk pada event source sesuai dengan skema \texttt{AuditEventDetails}, kemudian dienkripsi, ditandatangani, dan dikirimkan kepada ALS. ALS menyimpan setiap event sebagai \texttt{EncryptedSignedEvent} tanpa melakukan parsing terhadap isi payload (\textit{ciphertext}) pada saat penerimaan.

Hasil pengujian menunjukkan seluruh 13 jenis \textit{event} berhasil diterima dan disimpan. Seluruh 13 amplop event berhasil diterima dan disimpan sebagai \texttt{AuditRecord}. Jenis dan atribut masing-masing event tidak terlihat langsung pada berkas JSON-Lines karena berada dalam kondisi terenkripsi sebagai \textit{ciphertext}. Oleh karena itu, pengujian ini membuktikan kemampuan ALS menerima dan menyimpan 13 event terenkripsi yang dibangkitkan oleh event source, tanpa adanya parsing dari informasi event itu sendiri dari EV1-EV13 pada saat penerimaan.

\subsubsection{P-KF-03: Standardisasi Format \textit{Audit Record}}

Pengujian P-KF-03 memverifikasi KF-TA-03, yaitu kemampuan ALS memformat dan menstandardisasi \textit{audit record} ke dalam skema terstruktur sebelum disimpan. Skenario pengujian dilakukan dengan mengirimkan satu \textit{signed event} bertipe EV4 (\textit{Medical Record Access}) ke ALS, kemudian memeriksa struktur JSON dari \textit{audit record} yang dihasilkan.

\textit{Audit record} yang dihasilkan diverifikasi terhadap skema \texttt{AuditRecord} pada Kode~\ref{code:audit-record}. Pemeriksaan mencakup atribut rantai dan seluruh bidang \texttt{EncryptedSignedEvent}. Atribut umum \texttt{AuditEvent} beserta atribut tambahan EV4 hanya dapat diperiksa setelah \textit{ciphertext} didekripsi menggunakan kunci privat ALS.

Hasil pengujian menunjukkan \textit{audit record} yang dihasilkan memiliki
struktur yang sesuai dengan skema yang ditentukan. Seluruh bidang rantai
tersedia dan berformat benar, dan isi \textit{event} tersimpan sebagai
\texttt{EncryptedSignedEvent}. Setelah \textit{ciphertext} didekripsi dengan
kunci privat ALS, atribut tambahan EV4 berupa \texttt{access\_type},
\texttt{medical\_record\_id}, \texttt{capability\_id}, dan
\texttt{authorization\_token\_id} terserialisasi pada tingkat yang sama dengan
atribut umum berkat anotasi \texttt{\#[serde(flatten)]} pada
\texttt{AuditEvent}. 

\subsubsection{P-KF-04: Penyimpanan \textit{Audit Record} secara \textit{Append-Only}}

Pengujian P-KF-04 memverifikasi KF-TA-04, yaitu kemampuan ALS menyimpan \textit{audit record} pada media penyimpanan yang bersifat \textit{append-only}. Skenario pengujian dilakukan dengan mengirimkan lima \textit{event} secara berurutan ke ALS, kemudian memeriksa berkas log untuk memastikan setiap \textit{event} tersimpan sebagai baris baru tanpa menimpa baris sebelumnya.

Pemeriksaan dilakukan dengan menghitung jumlah baris pada berkas log sebelum dan sesudah pengiriman lima \textit{event}. Jumlah baris setelah pengiriman harus bertambah tepat lima baris dari jumlah sebelumnya. Selain itu, isi baris-baris lama diverifikasi untuk memastikan tidak ada perubahan.

Hasil pengujian menunjukkan jumlah baris bertambah tepat lima setelah pengiriman, dan isi baris-baris lama tetap tidak berubah. Penggunaan \texttt{OpenOptions::new().create(true).append(true)} pada \texttt{write\_audit\_record} terbukti merealisasikan karakteristik \textit{append-only} yang disyaratkan oleh KNF-TA-01.

\subsubsection{P-KF-05: Retensi Jangka Panjang melalui Rotasi Log, Pengarsipan IPFS, dan Publikasi Metadata IOTA}

Pengujian P-KF-05 memverifikasi KF-TA-05, yaitu kemampuan ALS menyediakan mekanisme retensi jangka panjang. Skenario pengujian dilakukan dengan mengonfigurasi \texttt{LOG\_ROTATION\_INTERVAL\_SECS} ke nilai kecil (60 detik) untuk mempercepat siklus rotasi, kemudian mengisi berkas log dengan sejumlah \textit{event} dan menunggu siklus rotasi berjalan.

Keberhasilan pengujian diverifikasi melalui tiga indikator: (1) berkas arsip \texttt{logs/uploading\_\{timestamp\}.log} terbentuk selama proses pengarsipan kemudian terhapus setelah selesai, (2) log konsol menampilkan pesan \texttt{[rotation] upload IPFS berhasil. CID: \{cid\}} dengan nilai CID yang valid, dan (3) log konsol menampilkan pesan \texttt{[iota] LogRecord dibuat. Object ID: \{id\} | TX: \{digest\}} yang menandakan metadata berhasil diterbitkan pada jaringan IOTA.

Hasil pengujian menunjukkan seluruh tiga indikator terpenuhi dalam lingkungan prototipe dan skenario uji yang ditetapkan. CID yang dihasilkan oleh IPFS dapat digunakan untuk mengambil kembali berkas arsip melalui \textit{gateway} IPFS, dan \textit{object ID} \texttt{LogRecord} pada IOTA dapat diverifikasi melalui IOTA Explorer pada jaringan Localnet.

\subsection{Pengujian End-to-End Pembangkitan, Pengiriman, dan Penyimpanan Audit
Record}
\label{subsec:pengujian-siklus-lengkap}
 
Pengujian end-to-end (P-INT-01) merupakan pengujian integrasi yang
memverifikasi keseluruhan \textit{pipeline} audit log dari pembangkitan
\textit{event} hingga verifikasi integritas rantai \textit{hash}, menggunakan
satu skenario konkret yang dapat direproduksi. Pengujian ini tidak menggantikan
P-KF atau P-KNF yang menguji komponen secara atomik, melainkan melengkapinya
dengan membuktikan bahwa seluruh komponen bekerja benar secara terpadu dalam
satu alur nyata.
 
Skenario yang digunakan adalah mengeksekusi transaksi ke IOTA yaitu update informasi rekam medis pasien oleh pasien sendiri dari \textit{Patient Client}, yang menghasilkan \textit{event}
jenis EV8 (\textit{IOTA Transaction}).
 
\subsubsection{Langkah 1: Event yang Terbentuk}
 
Ketika pasien melakukan sebuah operasi yang mengtrigger transaksi ke IOTA, \textit{Patient Client} membangkitkan
objek \texttt{Event} berikut dan membungkusnya ke dalam \texttt{AuditEvent}
sebelum dikirimkan ke ALS.
 
\begin{minted}[linenos, breaklines]{json}
{
  "source_component": "PatientClient",
  "source_timestamp": "2026-09-22T03:25:23.481414487Z",
  "event": {
    "actor_id": "0x24ec61172735ba2ca03c70ed6eb0a5301cd10380dc73cf2f41ad1c1916a9812d",
    "actor_type": "Patient",
    "target_object_type": "Transaction",
    "target_object": "tx://7sV8geo2CdtyhaPvn4UgktvkgRYRRQZttdKK7bYaud2r",
    "outcome": "Success",
    "action_type": "EXECUTE",
    "event_type": "EV8",
    "move_module": "patient",
    "move_function": "update_administrative_metadata"
  }
}
\end{minted}
 
Atribut \texttt{actor\_id} berisi alamat IOTA aktor yang juga digunakan untuk
menandatangani \textit{ciphertext} \texttt{EncryptedSignedEvent}. Atribut
\texttt{source\_timestamp} dicatat di sisi \textit{Patient Client} sebelum
enkripsi, sehingga nilainya mencerminkan waktu kejadian menurut jam komponen
pengirim dan tidak dapat diubah oleh ALS.
 
Sebelum dikirimkan, \texttt{AuditEvent} di atas dienkripsi menggunakan
AES-256-GCM dengan kunci sesi acak, kunci sesi dibungkus dengan kunci publik
P-256 ALS melalui ECIES, dan \textit{ciphertext} ditandatangani menggunakan
\texttt{IotaKeyPair} aktor. Hasilnya dikirimkan ke ALS dalam format
\texttt{EncryptedSignedEvent}.
 
\subsubsection{Langkah 2: Record yang Tersimpan}
 
Setelah ALS memverifikasi tanda tangan, mendekripsi \textit{payload}, dan
memproses \textit{event} melalui \texttt{AuditLogger}, satu baris JSON-Lines
berikut ditulis secara \textit{append} ke berkas \texttt{./logs/audit\_logs.log}.
Contoh \textit{record} kedua dalam rantai ditunjukkan di bawah (nilai
\texttt{prev\_record\_hash} mengacu pada \textit{record} sebelumnya).
 
\begin{minted}[linenos, breaklines]{json}
{
  "record_id": "01a0c726-f0e1-74e2-ae52-fd77e0d3d14b",
  "timestamp": "2026-09-22T03:26:53.153872401Z",
  "prev_record_hash": "e7c7954ff847965f1ae47010fe4c5e7adf8721d4a991c4253ec239ec899ae96c",
  "record_hash": "aeabafbc0e5b38d2669305448fcc8d90a83c68bf988c80a9e9e0e5f3d203d99e",
  "source_component": "PatientClient",
  "source_timestamp": "2026-09-22T03:25:23.481414487Z",
  "event": {
    "actor_id": "0x24ec61172735ba2ca03c70ed6eb0a5301cd10380dc73cf2f41ad1c1916a9812d",
    "actor_type": "Patient",
    "target_object_type": "Transaction",
    "target_object": "tx://7sV8geo2CdtyhaPvn4UgktvkgRYRRQZttdKK7bYaud2r",
    "outcome": "Success",
    "action_type": "EXECUTE",
    "event_type": "EV8",
    "move_module": "patient",
    "move_function": "update_administrative_metadata"
  }
}
\end{minted}
 
Selisih antara \texttt{source\_timestamp} (\texttt{03:25:23.4814Z}) dan
\texttt{timestamp} (\texttt{03:26:53.1538Z}) sebesar ±277 ms merupakan latensi
dari enkripsi, pengiriman jaringan, dan pemrosesan antrian internal ALS.
Kedua nilai diperlukan: \texttt{source\_timestamp} untuk merekonstruksi urutan
kejadian yang sesungguhnya, \texttt{timestamp} untuk menentukan urutan
pemrosesan di sisi ALS.
 
Nilai \texttt{record\_hash} dihitung oleh \texttt{Utils::calculate\_record\_hash}
sebagai SHA-256 dari gabungan \texttt{record\_id}, \texttt{timestamp},
\texttt{prev\_record\_hash}, dan seluruh isi \textit{event} yang diserialisasi
ke JSON. Proses ini dapat direproduksi: menjalankan fungsi yang sama terhadap
data yang sama selalu menghasilkan nilai yang identik.
 
\subsubsection{Langkah 3: Hubungan Antar-Record dan Rantai Hash}
 
Properti rantai \textit{hash} dapat diverifikasi dengan membandingkan nilai
\texttt{prev\_record\_hash} pada setiap \textit{record} dengan nilai
\texttt{record\_hash} pada \textit{record} sebelumnya. Untuk tiga \textit{record}
berurutan pada berkas log, hubungan yang harus terbentuk adalah:
 
\begin{center}
\begin{tabular}{|c|c|c|}
\hline
\textbf{N} & \textbf{\texttt{prev\_record\_hash}} & \textbf{\texttt{record\_hash}} \\
\hline
1 & \texttt{null} & \texttt{e7c795...} \\
\hline
2 & \texttt{e7c795...} & \texttt{aeabaf...} \\
\hline
\hline
\end{tabular}
\end{center}
 
Nilai \texttt{prev\_record\_hash} pada \textit{record} ke-$N$ harus sama persis
dengan \texttt{record\_hash} pada \textit{record} ke-$(N-1)$. Properti ini
diverifikasi melalui fungsi \texttt{verify\_chain} yang membaca seluruh baris
berkas log secara berurutan dan memeriksa konsistensi setiap pasangan
\textit{record} yang bersebelahan.
 
\subsubsection{Langkah 4: Verifikasi Integritas dan Deteksi Perusakan}
 
Verifikasi integritas dijalankan menggunakan \textit{unit test} \texttt{verify\_chain\_valid}
dan \texttt{verify\_chain\_tampered} dengan perintah:
 
\begin{minted}[breaklines]{bash}
cargo test verify_chain -- --nocapture
\end{minted}
 
\textbf{Skenario a --- Rantai valid.} Tiga \textit{record} dibuat secara
berurutan melalui \texttt{create\_audit\_record}. Fungsi \texttt{verify\_chain}
dipanggil terhadap ketiga \textit{record} tersebut. Hasil yang diharapkan adalah
seluruh \textit{record} dinyatakan valid dan fungsi mengembalikan \texttt{true}.
 
\textbf{Skenario b --- Perusakan pada posisi tengah.} Setelah rantai tiga
\textit{record} dibuat, nilai bidang \texttt{actor\_id} pada \textit{record}
kedua diubah secara langsung di memori. Fungsi \texttt{verify\_chain} dipanggil
kembali. Hasil yang diharapkan adalah \textit{record} kedua gagal verifikasi
karena nilai \texttt{record\_hash} yang tersimpan tidak lagi sesuai dengan
\textit{hash} yang dihitung ulang dari konten yang telah dimodifikasi. Selain
itu, \textit{record} ketiga juga gagal verifikasi karena nilai
\texttt{prev\_record\_hash}-nya tidak lagi cocok dengan \texttt{record\_hash}
\textit{record} kedua yang seharusnya.
 
Hasil pengujian yang terekam adalah: skenario a mengembalikan
\texttt{chain\_valid: true} untuk seluruh 3 \textit{record}, sedangkan skenario
b mengembalikan \texttt{chain\_valid: false} dengan pesan \texttt{record 2:
hash mismatch} dan \texttt{record 3: prev\_hash mismatch}. Hasil ini membuktikan bahwa modifikasi pada suatu \textit{record} terdeteksi
secara otomatis melalui efek domino pada \textit{record} sesudahnya.
 
\subsubsection{Langkah 5: Kesimpulan dari Satu Record}
 
Atribut \textit{event} pada \textit{record} EV4 tidak dapat dibaca langsung
dari berkas log karena tersimpan sebagai \textit{ciphertext}. Untuk keperluan
investigasi, pemegang kunci privat P-256 ALS mendekripsi \texttt{enc\_aes\_key}
untuk memperoleh kunci sesi AES, kemudian mendekripsi \texttt{ciphertext}
menjadi \texttt{AuditEvent}. Dari hasil dekripsi \textit{record} EV4 tersebut,
dapat disimpulkan:
 
\begin{itemize}
    \item \textbf{Siapa} yang mengakses: \texttt{actor\_id}
    \texttt{0xabc123...} dengan peran \texttt{MedicalPersonnel}.
    \item \textbf{Apa} yang diakses: RME dengan ID \texttt{RME-2025-00471}
    yang tersimpan di IPFS pada CID \texttt{bafybeig...}.
    \item \textbf{Kapan} akses terjadi: \texttt{2025-06-10T08:14:32.104Z}
    menurut jam komponen pengirim.
    \item \textbf{Dengan otorisasi apa}: \textit{capability}
    \texttt{cap:0xdef456...} yang dapat diverifikasi \textit{on-chain} di IOTA,
    dan JWT dengan ID \texttt{jwt:ey...}.
    \item \textbf{Hasilnya}: akses berhasil (\texttt{outcome: Success}).
\end{itemize}
 
Tanda tangan \texttt{IotaSignature} yang valid atas \textit{ciphertext}
membuktikan bahwa \textit{event} ini dilaporkan oleh pemegang kunci privat
IOTA yang bersesuaian dengan \texttt{iota\_address} pada \textit{record}.
Karena \texttt{signature} dan \texttt{iota\_address} ikut tersimpan pada
\textit{record}, pembuktian ini dapat diulang kapan saja terhadap
\textit{ciphertext} yang sama. Untuk mengatribusikan tindakan kepada
\texttt{actor\_id}, investigator mencocokkan \texttt{actor\_id} hasil dekripsi
dengan \texttt{iota\_address} penanda tangan. Apabila keduanya sama, pihak
tersebut tidak dapat menyangkal tindakan ini selama kunci privat IOTA-nya tidak
dikompromikan. Apabila keduanya berbeda, \textit{record} hanya membuktikan bahwa
pemilik \texttt{iota\_address} melaporkan tindakan tersebut atas nama
\texttt{actor\_id}, sehingga atribusi kepada \texttt{actor\_id} memerlukan
bukti pendukung lain.

\subsubsection{Langkah 6: Batas Kesimpulan}
 
Terdapat tiga batas yang perlu dinyatakan secara eksplisit agar kesimpulan
di atas tidak melampaui apa yang sesungguhnya dibuktikan oleh sistem.
 
\textbf{(a) Isi \textit{event} hanya dapat dibaca oleh pemegang kunci privat
ALS.} Berkas \texttt{audit\_logs.log} menyimpan \texttt{AuditRecord} dalam
format JSON-Lines, dengan atribut rantai (\texttt{record\_id},
\texttt{timestamp}, \texttt{prev\_record\_hash}, \texttt{record\_hash}) dalam
bentuk \textit{plaintext}, sedangkan isi \textit{event} tersimpan sebagai
\textit{ciphertext} AES-256-GCM. Hal yang sama berlaku untuk berkas arsip pada
IPFS, sehingga pihak yang mengetahui CID arsip tidak dapat membaca atribut
sensitif \textit{event}. Konsekuensinya, investigasi terhadap isi \textit{event}
bergantung pada ketersediaan kunci privat P-256 ALS; hilangnya kunci tersebut
menyebabkan isi seluruh \textit{event} tidak dapat dipulihkan, meskipun
integritas rantai tetap dapat diverifikasi. Selain itu, bidang
\texttt{iota\_address} pada amplop tersimpan dalam bentuk \textit{plaintext},
sehingga identitas penanda tangan setiap \textit{record} tetap dapat diketahui
tanpa dekripsi.
 
\textbf{(b) Fasilitas pencarian dan filter lintas-record belum tersedia.}
Implementasi saat ini hanya menyediakan \textit{endpoint} \texttt{GET /api/logs}
untuk membaca metadata rotasi \textit{on-chain}. Analisis lintas-\textit{record}
--- misalnya menemukan semua akses oleh aktor tertentu dalam rentang waktu
tertentu --- harus dilakukan secara manual terhadap berkas JSON-Lines atau
menggunakan alat bantu eksternal seperti \texttt{jq}. Fasilitas ini belum
terintegrasi ke dalam ALS.
 
\textbf{(c) Non-repudiation antar-sesi belum terjamin penuh.} Tanda tangan
menggunakan \texttt{IotaKeyPair} yang diambil dari \texttt{AppState} aktif.
Karena kunci IOTA terikat pada identitas aktor di \textit{ledger}, \textit{non-repudiation}
berlaku selama kunci privat tidak dikompromikan dan aktor masih memiliki kendali
atas kunci tersebut. Namun, apabila terjadi pergantian kunci IOTA aktor
antara waktu kejadian dan waktu investigasi, verifikasi tanda tangan terhadap
kunci publik saat ini dapat gagal meskipun \textit{event} tersebut sah.
Mekanisme rotasi kunci dengan riwayat kunci publik yang terekam berada di luar
cakupan implementasi saat ini.

\textbf{(d) \texttt{actor\_id} tidak diverifikasi oleh ALS pada saat
penerimaan.} ALS hanya memverifikasi bahwa tanda tangan dibuat oleh pemilik
\texttt{iota\_address}. Pelapor yang sah secara teknis dapat mengisi
\texttt{actor\_id} dengan identitas pengguna lain, dan \textit{event} tersebut
tetap diterima serta tercatat. Manipulasi ini tidak menghilangkan jejak pelapor,
karena \texttt{iota\_address} dan tanda tangannya tersimpan secara permanen
pada \textit{record} dan terlindungi oleh \textit{hash-chaining}, sehingga
ketidaksesuaian dapat terdeteksi dan dilacak ke pelapornya pada saat
investigasi. Namun, deteksi tersebut bergantung pada investigator yang
melakukan pencocokan secara eksplisit.

\textbf{(e) Integritas ujung rantai bergantung pada rotasi.}
\texttt{verify\_chain} memeriksa konsistensi antar \textit{record} yang ada,
tetapi tidak dapat mengetahui apakah \textit{record} terakhir telah dihapus.
Untuk \textit{record} yang telah dirotasi, kelengkapan dapat diverifikasi dengan
membandingkan \texttt{record\_hash} terakhir dan jumlah \textit{record} pada
berkas arsip terhadap \texttt{final\_record\_hash} dan \texttt{record\_count}
pada \texttt{LogRecord} di IOTA. Untuk \textit{record} yang masih berada pada
berkas log aktif, belum ada pembanding eksternal, sehingga penghapusan
\textit{record} terakhir oleh pihak yang memiliki akses ke berkas tersebut tidak
dapat dideteksi. Selain itu, karena setiap \textit{restart} ALS memulai rantai
baru, sebuah berkas log dapat memuat lebih dari satu segmen rantai, dan
\textit{record} dengan \texttt{prev\_record\_hash} bernilai \texttt{null} di
tengah berkas harus diinterpretasikan sebagai titik \textit{restart}.
Penghapusan \textit{record} terakhir sebelum \textit{restart} tidak dapat
dibedakan dari akhir segmen yang sah.
 
Dengan demikian, P-INT-01 membuktikan bahwa mekanisme inti sistem berfungsi
benar: \textit{record} terbentuk sesuai skema, rantai \textit{hash} valid,
perusakan terdeteksi, dan dari satu \textit{record} sudah dapat disimpulkan
bukti minimal yang relevan untuk keperluan akuntabilitas. Klaim ini tidak
melampaui batas ketiga yang telah disebutkan di atas.

\subsection{Pengujian Kebutuhan Nonfungsional}
\label{subsec:pengujian-nonfungsional}

Pengujian menggunakan fungsi verifikasi tanda tangan yang merepresentasikan fungsi pada handler ALS. Proses dekripsi dilakukan melalui fungsi pembantu terpisah untuk merepresentasikan pembacaan audit record oleh pihak yang memiliki kunci privat ALS. \textit{Event} uji dibentuk dengan fungsi pembantu \texttt{seal\_event} yang
meniru langkah 1 hingga 5 pada \texttt{ALSClient::build\_and\_send}: serialisasi
\texttt{AuditEvent}, enkripsi AES-256-GCM, pembungkusan kunci sesi dengan ECIES
P-256, serta penandatanganan \textit{ciphertext} dengan \texttt{IotaKeyPair}.
Pasangan kunci aktor dan pasangan kunci P-256 ALS dibangkitkan secara acak pada
setiap pengujian.

Pengujian P-KNF-06 memverifikasi bahwa \textit{event} yang dibentuk secara sah
diterima dan isinya utuh setelah didekripsi. Pengujian P-KNF-07 memverifikasi
integritas pada tiga lapisan: modifikasi \textit{ciphertext} harus ditolak oleh
verifikasi \texttt{IotaSignature} (P-KNF-07a), modifikasi \textit{nonce} yang
tidak tercakup tanda tangan harus ditolak oleh \textit{authentication tag}
AES-GCM (P-KNF-07b), dan kunci sesi yang dibungkus untuk kunci publik lain harus
gagal dibuka oleh ALS (P-KNF-07c). Pengujian P-KNF-08 memverifikasi keaslian
pengirim: tanda tangan dari kunci lain yang mengaku sebagai alamat aktor sah
(P-KNF-08a) maupun penggantian \texttt{iota\_address} pada amplop (P-KNF-08b)
harus ditolak oleh \texttt{verify\_secure}. Pengujian P-KNF-09 memverifikasi kerahasiaan payload pada penyimpanan, yaitu bahwa atribut sensitif event tidak tersimpan dalam bentuk plaintext dan hanya dapat dipulihkan menggunakan kunci privat P-256 ALS.

\subsubsection{Pengujian KNF-TA-01: Integritas \textit{Audit Record} yang \textit{Tamper-Proof}}

Pengujian KNF-TA-01 memastikan bahwa \textit{audit record} yang telah tersimpan bersifat \textit{tamper-evident} dan perusakan dapat terdeteksi melalui mekanisme \textit{hash-chaining}. Pengujian ini menguji mitigasi terhadap ancaman T7 dan T8 (kategori \textit{Repudiation}) yang berkaitan dengan kemampuan tenaga medis dan admin fasyankes menyangkal tindakan yang telah dilakukan.

Pengujian pertama, yaitu P-KNF-01, dijalankan melalui unit testing pada fungsi \texttt{Utils::verify\_record}. Skenario pengujian menciptakan sebuah \texttt{AuditRecord} yang valid menggunakan \texttt{create\_audit\_record}, kemudian memverifikasi \textit{record} tersebut dengan \texttt{verify\_record}. Hasil verifikasi harus mengembalikan \texttt{true}. Selanjutnya, satu karakter pada bidang \texttt{source\_component} dalam \textit{record} dimodifikasi secara langsung, dan \texttt{verify\_record} dipanggil kembali. Hasil verifikasi harus mengembalikan \texttt{false} karena nilai \texttt{record\_hash} yang tersimpan tidak lagi sesuai dengan \textit{hash} yang dihitung ulang dari konten \textit{record} yang telah dimodifikasi.

Pengujian kedua, yaitu P-KNF-02, dijalankan melalui unit testing untuk mensimulasikan perusakan pada posisi tengah rantai. Tiga \texttt{AuditRecord} dibuat secara berurutan sehingga membentuk rantai (setiap \textit{record} menyertakan \texttt{prev\_record\_hash} dari \textit{record} sebelumnya). \textit{Record} kedua kemudian dimodifikasi, dan seluruh tiga \textit{record} diverifikasi. \textit{Record} pertama harus tetap valid, sedangkan \textit{record} kedua dan ketiga harus gagal verifikasi karena efek domino perusakan rantai.

Hasil pengujian menunjukkan bahwa modifikasi \textit{record} di tengah segmen rantai yang diuji dapat dideteksi melalui ketidaksesuaian \texttt{record\_hash} dan \texttt{prev\_record\_hash}. Pengujian ini tidak membuktikan deteksi penghapusan record terakhir pada log aktif maupun kontinuitas rantai setelah ALS dimulai ulang. Implementasi \texttt{calculate\_record\_hash} yang menggabungkan
\texttt{record\_id}, \texttt{timestamp}, \texttt{prev\_record\_hash}, dan isi
\textit{event} terbukti memberikan sifat \textit{tamper-evident} terhadap
modifikasi pada berkas log lokal. Pengujian ini tidak mencakup penghapusan
\textit{record} terakhir, yang tidak dapat dideteksi oleh \texttt{verify\_chain}
tanpa pembanding \textit{hash} ujung rantai dari luar berkas log.

\subsubsection{Pengujian KNF-TA-02: Interoperabilitas dan Keterbacaan Format Data Log}

Pengujian KNF-TA-02 memastikan format data log yang dihasilkan konsisten, terstruktur, dan mudah dibaca oleh komponen sistem. Pengujian ini berkaitan dengan kemampuan format JSON-Lines untuk diproses lintas komponen tanpa konversi tambahan.

Pengujian P-KNF-03 dijalankan melalui unit testing yang membaca kembali berkas log yang telah terisi beberapa \textit{audit record} dan melakukan deserialisasi setiap baris JSON menggunakan \texttt{serde\_json::from\_str::<AuditRecord>}. Seluruh baris harus berhasil dideserializasi tanpa \textit{error} menjadi objek \texttt{AuditRecord} yang valid, dan atribut-atributnya harus cocok dengan nilai yang dikirimkan saat pengiriman \textit{event}.

Hasil pengujian menunjukkan seluruh baris berkas log berhasil dideserializasi. Penggunaan \texttt{\#[serde(flatten)]} pada \texttt{AuditEvent} dan \texttt{AuditEventDetails} terbukti menghasilkan format JSON yang datar (\textit{flat}) dan mudah diproses, tanpa hierarki objek yang bersarang secara berlebihan.

\subsubsection{Pengujian KNF-TA-03: Ketersediaan Data \textit{Audit Log} dalam Jangka Panjang}

Pengujian KNF-TA-03 memastikan ketersediaan dan keteraksesan data \textit{audit log} dalam jangka panjang melalui kombinasi \textit{IPFS pinning} dan penyimpanan metadata \textit{on-chain} pada IOTA.

Pengujian P-KNF-04 dijalankan dengan mengunggah sebuah berkas arsip log ke IPFS melalui fungsi \texttt{Utils::add\_file\_to\_ipfs}, kemudian mengambil kembali berkas tersebut menggunakan CID yang diperoleh melalui \textit{endpoint} \texttt{GET \{IPFS\_GATEWAY\_BASE\_URL\}/ipfs/\{cid\}}. Konten yang diambil kembali harus identik dengan konten berkas yang diunggah, diverifikasi dengan membandingkan nilai \textit{hash} SHA-256 dari keduanya.

Pengujian P-KNF-05 dijalankan dengan memverifikasi bahwa objek \texttt{LogRecord} yang diterbitkan ke IOTA melalui \texttt{IotaLogClient::publish\_metadata}. Pemeriksaan dilakukan melalui \textit{endpoint} \texttt{GET /object/\{object\_id\}} pada IOTA Localnet RPC. Respons harus memuat field \texttt{"owner": "Immutable"} yang menandakan objek tidak dapat dimodifikasi oleh siapa pun setelah transaksi pembuatannya final.

Hasil pengujian P-KNF-04 menunjukkan berkas arsip yang diambil kembali dari IPFS memiliki nilai \textit{hash} SHA-256 yang identik dengan berkas aslinya, membuktikan integritas data selama penyimpanan dan pengambilan. Hasil pengujian P-KNF-05 menunjukkan objek \texttt{LogRecord} pada IOTA Localnet berstatus \textit{immutable}, membuktikan bahwa Move \texttt{als::audit\_log} berhasil menegakkan sifat \textit{append-only} pada level \textit{on-chain}.

\subsubsection{Pengujian KNF-TA-04: Keaslian, Integritas, dan Kerahasiaan \textit{Event}}
 
Pengujian KNF-TA-04 memastikan bahwa jalur kriptografi yang digunakan pada
implementasi, yaitu verifikasi \texttt{IotaSignature} pada pola Encrypt-then-Sign mencegah \textit{ciphertext} yang dimodifikasi atau ditandatangani oleh kunci yang tidak sesuai diterima sebagai event yang sah. Bidang amplop yang tidak tercakup tanda tangan, seperti \texttt{nonce} dan \texttt{enc\_aes\_key}, tetap tersimpan sebagai bagian \texttt{AuditRecord} dan dilindungi oleh \texttt{record\_hash}, tetapi karusakan bidang tersebut baru dapat diketahui ketika proses dekripsi dilakukan. Pengujian ini menguji mitigasi
terhadap ancaman T6 (pasien menyangkal pemberian akses) dan T16 (PRE Server
menyangkal pengajuan \textit{sponsorship}), karena \textit{event} yang masuk ke
ALS harus dapat dibuktikan berasal dari pemegang kunci IOTA yang sah.
 
Seluruh pengujian dijalankan melalui \textit{unit testing} terhadap fungsi yang
dipanggil pada alur produksi, yaitu \texttt{Utils::verify\_event\_signature}
pada \textit{handler} \texttt{handle\_event} dan \texttt{create\_audit\_record}
pada \texttt{AuditLogger}. \textit{Event} uji dibentuk dengan fungsi pembantu
\texttt{seal\_event} yang meniru langkah 1 hingga 5 pada
\texttt{ALSClient::build\_and\_send}: serialisasi \texttt{AuditEvent}, enkripsi
AES-256-GCM dengan kunci sesi acak, pembungkusan kunci sesi dengan ECIES P-256
terhadap kunci publik ALS, serta penandatanganan \textit{ciphertext} dengan
\texttt{IotaKeyPair} pengirim menggunakan \texttt{IotaSignature}. Pasangan kunci
IOTA pengirim dan pasangan kunci P-256 ALS dibangkitkan secara acak pada setiap
pengujian sehingga tidak ada kunci yang ditanam pada kode uji. Skenario
pengujian dirangkum pada Tabel~\ref{tab:skenario-knf-ta-04}.
 
\begin{table}[htbp]
\centering
\caption{Skenario Pengujian KNF-TA-04}
\label{tab:skenario-knf-ta-04}
\renewcommand{\arraystretch}{1.3}
\begin{tabular}{|p{1.7cm}|p{6.3cm}|p{4.5cm}|}
\hline
\textbf{ID} & \textbf{Skenario} & \textbf{Hasil yang Diharapkan} \\
\hline
P-KNF-06 & \texttt{EncryptedSignedEvent} yang dibentuk secara sah diverifikasi oleh ALS. & Diterima. \\
\hline
P-KNF-07a & Satu \textit{byte} \texttt{ciphertext} diubah setelah penandatanganan. & Ditolak pada verifikasi tanda tangan. \\
\hline
P-KNF-07b & Modifikasi \texttt{nonce} tidak terdeteksi pada tahap verifikasi tanda tangan karena \texttt{nonce} tidak termasuk data yang ditandatangani. Event tersebut dapat diterima dan tersimpan, tetapi proses dekripsi kemudian gagal akibat verifikasi authentication tag AES-GCM. Dengan demikian, skenario ini membuktikan deteksi saat pembacaan, bukan penolakan saat penerimaan. \\
\hline
P-KNF-08a & \textit{Ciphertext} ditandatangani kunci lain, tetapi \texttt{iota\_address} diisi alamat pengirim sah. & Ditolak pada verifikasi tanda tangan. \\
\hline
P-KNF-08b & \texttt{iota\_address} pada amplop diganti dengan alamat lain. & Ditolak pada verifikasi tanda tangan. \\
\hline
P-KNF-09a & \textit{Audit record} yang dibentuk dari \textit{event} EV4 diserialisasi ke JSON-Lines. & Tidak memuat atribut \textit{event} dalam bentuk \textit{plaintext}. \\
\hline
P-KNF-09b & Isi \textit{event} pada \textit{record} dibuka dengan kunci privat ALS dan dengan kunci privat P-256 lain. & Hanya kunci privat ALS yang berhasil; isi identik dengan \textit{event} asli. \\
\hline
\end{tabular}
\end{table}
 
Pengujian P-KNF-06 memverifikasi bahwa \textit{event} yang sah tidak ditolak
oleh ALS. Pengujian P-KNF-07 memverifikasi integritas. Pada P-KNF-07a,
modifikasi \textit{ciphertext} harus terdeteksi oleh \texttt{IotaSignature}
karena \textit{ciphertext} merupakan pesan yang ditandatangani. Pada P-KNF-07b,
bidang \texttt{nonce} tidak tercakup tanda tangan sehingga modifikasinya tidak
terdeteksi pada saat penerimaan; perlindungan terhadap bidang ini berasal dari
\textit{authentication tag} AES-GCM, yang menyebabkan dekripsi gagal pada saat
pembacaan log. Pengujian P-KNF-08 memverifikasi keaslian pengirim, yaitu bahwa
tanda tangan hanya diterima apabila dibuat oleh kunci privat yang bersesuaian
dengan \texttt{iota\_address} pada amplop. Pengujian P-KNF-09 memverifikasi
kerahasiaan pada penyimpanan: \textit{record} pada berkas log tidak memuat isi
\textit{event} yang dapat dibaca, dan isi tersebut hanya dapat dipulihkan oleh
pemegang kunci privat P-256 ALS.

\section{Evaluasi}
\label{sec:evaluasi}

Evaluasi dilakukan dengan membandingkan hasil implementasi dan pengujian terhadap kebutuhan sistem. Hasil evaluasi digunakan untuk menilai apakah pendekatan \textit{Centralized Audit Collector} dengan \textit{signed payload} Ed25519, \textit{hash-chaining}, pengarsipan IPFS, dan publikasi metadata \textit{on-chain} pada IOTA mampu menjawab keterbatasan mekanisme pencatatan pada DecMed awal yang telah diidentifikasi pada Subbab~\ref{analisis-decmed}.

Berdasarkan hasil pengujian, kebutuhan fungsional dan KNF-TA-01 sampai KNF-TA-03 dinyatakan terpenuhi dalam batas skenario prototipe, sedangkan KNF-TA-04 terpenuhi sebagian karena belum terdapat jaminan durabilitas end-to-end setelah ALS memberikan respons berhasil.

\subsection{Evaluasi terhadap Kebutuhan Sistem}
\label{subsec:evaluasi-kebutuhan}

Berdasarkan kebutuhan sistem pada Subbab~\ref{subsec:impl-skema-tipe-data} hingga~\ref{subsec:rancangan-publikasi-iota}, implementasi dievaluasi terhadap kebutuhan fungsional dan nonfungsional yang menjadi dasar pengembangan ALS. Setiap kebutuhan diperiksa berdasarkan komponen yang merealisasikannya dan skenario pengujian yang memverifikasi perilakunya. Pemetaan ini menunjukkan keterhubungan antara kebutuhan, implementasi, dan pengujian.

Ringkasan pemetaan kebutuhan sistem terhadap implementasi dan pengujian disajikan pada Tabel~\ref{tab:evaluasi-kebutuhan}.

\begin{xltabular}{\textwidth}{|p{1.8cm}|p{3.5cm}|p{3.5cm}|p{3.2cm}|}
\caption{Evaluasi Pemenuhan Kebutuhan Sistem \textit{Audit Log}}
\label{tab:evaluasi-kebutuhan}\\
\hline
\textbf{ID Kebutuhan} & \textbf{Komponen Implementasi} & \textbf{Skenario Pengujian} & \textbf{Status} \\
\hline
\endfirsthead
\hline
\textbf{ID Kebutuhan} & \textbf{Komponen Implementasi} & \textbf{Skenario Pengujian} & \textbf{Status} \\
\hline
\endhead
\hline
\endfoot

KF-TA-01 & \textit{Endpoint} \texttt{POST /api/events}, \texttt{AuditLogger}, antrian \texttt{tokio::sync::mpsc} & P-KF-01 & Terpenuhi \\
\hline
KF-TA-02 & \texttt{AuditEventDetails} \textit{tagged enum} (EV1--EV13) & P-KF-02 & Terpenuhi \\
\hline
KF-TA-03 & \texttt{AuditRecord} dengan \texttt{\#[serde(flatten)]} & P-KF-03 & Terpenuhi \\
\hline
KF-TA-04 & \texttt{write\_audit\_record} dengan mode \textit{append} & P-KF-04 & Terpenuhi \\
\hline
KF-TA-05 & Log Rotation Worker, \texttt{add\_file\_to\_ipfs}, \texttt{publish\_metadata}, \texttt{freeze\_object} & P-KF-05, P-KNF-04, P-KNF-05 & Terpenuhi \\
\hline
KNF-TA-01 & \texttt{create\_audit\_rec- ord} (\textit{hash-chaining}), \texttt{verify\_record}, \textit{ object on-chain} & P-KNF-01, P-KNF-02, P-KNF-05 & Terpenuhi \\
\hline
KNF-TA-02 & Format JSON-Lines, \texttt{\#[serde(flatten)]}, skema \texttt{AuditRecord} & P-KNF-03 & Terpenuhi \\
\hline
KNF-TA-03 & \textit{IPFS pinning}, \texttt{LogRecord} sebagai \textit{object immutable} di IOTA & P-KNF-04, P-KNF-05 & Terpenuhi \\
\hline
KNF-TA-04 & Verifikasi \texttt{IotaSignature} pada \texttt{verify\_event\_signature}, \texttt{AlsQueue} dan \textit{retry worker} pada \textit{event source} & P-KNF-06, P-KNF-07, P-KNF-08, P-KNF-09 & Terpenuhi sebagian \\
\hline
\end{xltabular}

Tabel~\ref{tab:evaluasi-kebutuhan} menunjukkan bahwa seluruh kebutuhan
fungsional (KF-TA-01 hingga KF-TA-05) serta kebutuhan nonfungsional KNF-TA-01
hingga KNF-TA-03 terpenuhi berdasarkan skenario pengujian yang dijalankan.
KNF-TA-04 dinyatakan terpenuhi sebagian. Verifikasi tanda tangan digital
terbukti mencegah \textit{event} palsu atau termanipulasi masuk ke rantai
\textit{audit record}, dan \texttt{AlsQueue} mengurangi risiko \textit{event}
tidak terkirim akibat kegagalan sementara. Namun, implementasi belum menjamin
bahwa setiap \textit{event} yang telah dikonfirmasi oleh ALS pasti tersimpan
pada berkas log, sebagaimana diuraikan pada
Subbab~\ref{subsec:keterbatasan-implementasi}.

Kebutuhan fungsional diverifikasi melalui skenario P-KF-01 sampai P-KF-05, sedangkan kebutuhan nonfungsional diverifikasi melalui pengujian P-KNF-01 sampai P-KNF-08.

Evaluasi terhadap masalah-masalah yang diidentifikasi pada Tabel~\ref{tab:identifikasi-masalah} menghasilkan temuan sebagai berikut. Masalah M-TA-01 (belum adanya sistem \textit{audit log} terintegrasi) diselesaikan melalui penambahan ALS sebagai \textit{centralized audit collector} dengan titik masuk tunggal \texttt{POST /api/events}, diverifikasi oleh P-KF-01. Masalah M-TA-02 (cakupan \textit{event} yang sempit) diselesaikan melalui katalog 13 jenis \textit{event} (EV1--EV13) yang diturunkan dari analisis ancaman STRIDE, diverifikasi oleh P-KF-02. Masalah M-TA-03 (\textit{audit record} yang masih dapat dimodifikasi) diselesaikan melalui mekanisme \textit{hash-chaining} pada berkas log lokal dan penyimpanan pada \textit{on-chain} melalui \textit{smart contract} Move, diverifikasi oleh P-KNF-01, P-KNF-02, dan P-KNF-05. Masalah M-TA-04 (format \textit{audit record} yang tidak terstandarisasi) diselesaikan melalui skema \texttt{AuditRecord} yang seragam dengan format JSON-Lines, diverifikasi oleh P-KNF-03. Masalah M-TA-05 (kerentanan terhadap \textit{data pruning}) diselesaikan melalui strategi persistensi ganda berbasis \textit{IPFS pinning} dan penyimpanan \textit{state object} Move, diverifikasi oleh P-KNF-04 dan P-KNF-05.

\subsection{Keterbatasan Implementasi}
\label{subsec:keterbatasan-implementasi}

Implementasi pada Tugas Akhir ini memiliki beberapa keterbatasan sebagai berikut.

\begin{enumerate}
    \item Implementasi dilakukan dalam lingkungan prototipe yang terdiri atas komputer lokal, VPS, dan jaringan IOTA Localnet. Seluruh komponen ALS dan PRE Server dijalankan pada komputer lokal dengan sistem operasi Ubuntu 24.04.4 di lingkungan WSL2. VPS menjalankan IPFS Cluster dan IOTA Gas Station. \textit{Smart contract} \texttt{als::audit\_log} diunggah ke jaringan IOTA Localnet, bukan jaringan IOTA Testnet maupun Mainnet, sehingga perilaku sistem pada jaringan publik yang sesungguhnya belum dapat diverifikasi.
    
    \item Mekanisme \textit{IPFS pinning} yang diimplementasikan bergantung pada konfigurasi IPFS Cluster yang diwarisi dari implementasi DecMed sebelumnya dengan tiga \textit{peer}. Ketahanan \textit{pinning} terhadap kegagalan \textit{node} pada kluster berskala produksi belum diuji.
    \item Nilai \texttt{LOG\_ROTATION\_INTERVAL\_SECS} yang digunakan dalam pengujian (60 detik) berbeda dari nilai yang direkomendasikan untuk produksi (300 detik sesuai konstanta). Implikasi rotasi yang terlalu sering terhadap konsumsi \textit{gas} IOTA dan kuota transaksi Gas Station belum dievaluasi secara menyeluruh.
    
    \item Kunci privat Ed25519 yang digunakan oleh setiap \textit{event source} untuk menandatangani \textit{audit event} dibangkitkan secara acak pada saat komponen diinisialisasi dan tidak dipersistensikan. Apabila komponen di-\textit{restart}, kunci baru akan dibangkitkan sehingga tidak ada mekanisme untuk menghubungkan \textit{event} dari sesi sebelumnya dengan sesi baru secara kriptografis. Pengelolaan kunci jangka panjang untuk keperluan \textit{non-repudiation} yang sesungguhnya memerlukan mekanisme penyimpanan kunci yang aman dan persisten.
    
    \item Pengujian verifikasi integritas rantai \textit{audit record} dilakukan pada skala kecil (puluhan \textit{record}). Performa verifikasi rantai pada volume berkas log yang besar (ribuan hingga jutaan \textit{record}) belum dievaluasi.
    
    \item Mekanisme pengiriman \textit{event} mengurangi risiko kehilangan \textit{event}, tetapi tidak memberikan jaminan ketercapaian \textit{end-to-end}. Terdapat tiga kondisi yang dapat menyebabkan \textit{event} tidak tercatat. Pertama, ALS mengirim respons berhasil setelah \textit{event} masuk ke antrian \texttt{tokio::sync::mpsc} di memori, sebelum \textit{record} ditulis ke berkas log, sedangkan \textit{event source} menghapus entri dari \texttt{AlsQueue} segera setelah respons tersebut diterima; kegagalan ALS di antara kedua titik ini menyebabkan \textit{event} hilang. Kedua, \textit{event} dengan tanda tangan tidak valid ditolak dengan status HTTP 200, sehingga pengirim memperlakukannya sebagai pengiriman berhasil dan menghapusnya dari antrian tanpa jejak penolakan. Ketiga, entri yang melampaui \texttt{MAX\_ATTEMPTS\_BEFORE\_SKIP} tidak lagi dikirim ulang secara otomatis dan memerlukan penanganan manual oleh operator. Mitigasi yang dapat diterapkan pada pengembangan selanjutnya adalah mengirim respons berhasil hanya setelah \textit{record} tersimpan ke disk, menggunakan status HTTP yang membedakan penolakan dari keberhasilan, serta menyediakan mekanisme pemantauan entri antrian yang gagal terkirim.
    
    \item ALS tidak menegakkan kesesuaian antara \texttt{actor\_id} di dalam
    \textit{event} dan \texttt{iota\_address} penanda tangan, karena isi
    \textit{event} tidak didekripsi pada saat penerimaan. Akibatnya, pelapor
    yang sah berpotensi mencatat \textit{event} dengan \texttt{actor\_id} milik
    pengguna lain. Risiko ini dimitigasi sebagian karena identitas pelapor
    tetap tersimpan dan terverifikasi pada setiap \textit{record}, sehingga
    atribusi yang tidak sesuai dapat terdeteksi dan ditelusuri ke pelapornya
    pada saat investigasi. Penegakan otomatis dapat ditambahkan pada
    pengembangan selanjutnya, misalnya dengan menolak \textit{event} yang
    \texttt{actor\_id}-nya berbeda dari \texttt{iota\_address} pada saat
    dekripsi, atau dengan menambahkan daftar pelapor tepercaya yang diizinkan
    mencatat tindakan atas nama aktor lain.

    \item Nilai \texttt{prev\_record\_hash} hanya dipelihara di memori dan tidak dipulihkan setelah ALS dimulai ulang,
    sehingga setiap \textit{restart} memutus keterikatan antara \textit{record}
    sebelum dan sesudah \textit{restart}. Selain itu, \textit{hash chain} tidak
    mendeteksi penghapusan \textit{record} terakhir pada berkas log aktif sebelum nilai ujung rantai dipublikasikan ke IOTA melalui rotasi. Rentang kerentanan ini dibatasi oleh \texttt{LOG\_ROTATION\_INTERVAL\_SECS}: semakin
    pendek interval rotasi, semakin sedikit \textit{record} yang belum memiliki
    \textit{anchor} eksternal. Mitigasi yang dapat diterapkan pada pengembangan
    selanjutnya meliputi pemulihan \texttt{prev\_record\_hash} dari \textit{record} terakhir pada berkas log aktif atau dari \texttt{final\_record\_hash} pada IOTA saat inisialisasi, serta publikasi \textit{hash} ujung rantai secara lebih sering, misalnya setiap sejumlah \textit{record} tertentu, tanpa menunggu rotasi berkas. Selain itu, nilai \texttt{prev\_tx\_digest} yang menghubungkan metadata antarrotasi juga dipelihara di memori dan belum dipulihkan dari IOTA setelah restart. Oleh karena itu, restart ALS berpotensi memutus bukan hanya rantai antarrecord, tetapi juga rantai metadata antarrotasi.

    \item Implementasi menggunakan pasangan kunci P-256 ALS untuk membuka kunci sesi seluruh audit event, tetapi belum menyediakan mekanisme identifikasi versi kunci, rotasi kunci, backup, escrow, maupun pemulihan kunci. Kehilangan kunci privat ALS menyebabkan event yang dibungkus untuk kunci tersebut tidak dapat didekripsi, sedangkan kompromi kunci dapat membuka seluruh arsip terkait. Pengembangan selanjutnya perlu menambahkan \texttt{key\_id} atau \texttt{key\_version} pada amplop event serta kebijakan penyimpanan, rotasi, dan retensi kunci lama.
 
\end{enumerate}

Keterbatasan tersebut menunjukkan bahwa implementasi masih berada pada tingkat pembuktian konsep (\textit{proof of concept}). Pengembangan berikutnya perlu menguji sistem pada jaringan IOTA Testnet, memperkuat mekanisme pengelolaan kunci Ed25519 untuk mendukung \textit{non-repudiation} jangka panjang, serta mengevaluasi skalabilitas \textit{hash-chain} verification pada volume log produksi.